\documentclass[11pt]{article}
\usepackage{amsmath,amssymb,booktabs,geometry,microtype}
\geometry{margin=1in}

\newcommand{\css}{\mathrm{css}}
\newcommand{\fin}{\mathrm{fin}}

\title{Infinite-Horizon Tail: Current Formulation and the Effect of Discounting}
\date{}

\begin{document}
\maketitle

\section{Notation}

The controller uses a finite horizon $[0,T_f]$ followed by a compactified
infinite tail $[T_f,\infty)$ mapped onto $\tau\in[0,1)$ by
%
\begin{equation}
  t \;=\; T_f + L\,\operatorname{atanh}\tau,
  \qquad
  \gamma \;=\; \frac{1}{L},
  \qquad
  \frac{\mathrm{d}t}{\mathrm{d}\tau} \;=\; \frac{1}{\gamma\bigl(1-\tau^{2}\bigr)} .
  \label{eq:transform}
\end{equation}
%
Thus $\tau=1$ corresponds to $t=\infty$. The tail is discretised with
$n_{\mathrm{fe}}=5$ elements and $n_{\mathrm{cp}}=3$ collocation points per
element. Let $\mathcal{K}$ denote the set of quadrature nodes.

\begin{center}
\begin{tabular}{ll}
\toprule
Symbol & Meaning \\
\midrule
$\rho_{p,v}(\tau)$ & pipe density (differential state) \\
$p_{p,v}(\tau)$    & interior pipe pressure \texttt{interm\_p} (algebraic) \\
$\beta_s(\tau)$    & compressor pressure ratio (control) \\
$P_s(\tau)$        & compressor power \\
$w_v,\;c_v$        & cell inlet flow and local consumption \\
$(\cdot)^{\css}$   & cyclic-steady-state reference \\
\bottomrule
\end{tabular}
\end{center}

\section{Current Formulation}

\subsection{Tail quadrature}

Every tail integral is evaluated as an explicit weighted sum,
%
\begin{equation}
  \Phi_\bullet \;=\; \sum_{k\in\mathcal{K}} \frac{w_k}{\gamma}\,
                     \ell_\bullet(\tau_k),
  \qquad
  \boxed{\;\tau=1 \notin \mathcal{K}\;}
  \label{eq:quad}
\end{equation}
%
with the transformation Jacobian folded into the weights,
%
\begin{equation}
  w_k \;=\;
  \begin{cases}
    \displaystyle\int_{\mathrm{elem}} L_k(\tau)\,\frac{\mathrm{d}\tau}{1-\tau^{2}},
      & \text{\textsc{lagrange-radau}},\\[2.2ex]
    \displaystyle\frac{h_i\,\omega_k}{1-\tau_k^{2}},
      & \text{\textsc{lagrange-legendre}} .
  \end{cases}
  \label{eq:weights}
\end{equation}
%
The node $\tau=1$ is excluded because the Jacobian $1/(1-\tau^{2})$ is singular
there. On a Radau mesh $\tau=1$ \emph{is} a collocation point, and the
remaining Lagrange bases satisfy $L_k(1)=0$; that root cancels the pole
analytically,
%
\begin{equation}
  \frac{L_k(\tau)}{1-\tau^{2}}
  \;=\;\frac{(1-\tau)\,q_k(\tau)}{(1-\tau)(1+\tau)}
  \;=\;\frac{q_k(\tau)}{1+\tau},
\end{equation}
%
so all weights in \eqref{eq:weights} are finite.

\subsection{Stage costs}

Three stage costs are integrated over the tail, and they are \emph{not} of the
same type:
%
\begin{align}
  \ell_P(\tau)      &= \sum_{s}\frac{P_s(\tau)}{1000}
                       &&\text{absolute, economic}, \label{eq:lP}\\
  \ell_V(\tau)      &= \sum_{(p,v)}\bigl(p_{p,v}(\tau)-p^{\css}_{p,v}(\tau)\bigr)^{2}
                       &&\text{tracking}, \\
  \ell_\beta(\tau)  &= \sum_{s}\bigl(\beta_s(\tau)-\beta^{\css}_s(\tau)\bigr)^{2}
                       &&\text{tracking}.
\end{align}

\subsection{Objective}

\begin{equation}
  J \;=\;
  \underbrace{\sum_{t<t_f}\sum_{s}\frac{P_{s,t}\,\Delta t}{1000}}_{\text{finite horizon}}
  \;+\;
  \underbrace{3600\,\sigma\,\Phi_P}_{\text{tail}}
  \;+\;
  \underbrace{\mu\sum_{p,v}\bigl(s^{+}_{p,v}+s^{-}_{p,v}\bigr)}_{\text{terminal }L_1}
  \label{eq:obj}
\end{equation}
%
with $\Delta t=3600\,\mathrm{s}$ and $\sigma=1$ (no tail overestimation
factor). The Lyapunov function and its descent inequality are imposed
separately,
%
\begin{equation}
  V \;=\; \frac{p_{\fin}}{f_p} + \frac{\Phi_V}{\inf_p}
        + \frac{\beta_{\fin}}{f_\beta} + \frac{\Phi_\beta}{\inf_\beta},
  \qquad
  V_k \;\le\; V_{k-1} - \delta\,\ell^{\mathrm{plant}}_{k-1}.
  \label{eq:lyap}
\end{equation}

\subsection{Terminal condition ($L_1$-soft)}

\begin{equation}
  p_{p,v}(\tau{=}1) \;=\; p^{\css}_{p,v} + s^{+}_{p,v} - s^{-}_{p,v},
  \qquad s^{\pm}_{p,v}\ge 0,
  \label{eq:terminal}
\end{equation}
%
penalised in \eqref{eq:obj} by $\mu$. The $L_1$ form is an \emph{exact}
penalty: above a finite threshold its solution coincides with the
hard-constrained one wherever that is feasible, and it concentrates any
residual violation on the rows that genuinely cannot be met.

\subsection{Equations retained at $\tau=1$}

On a Radau mesh $\tau=1$ is a collocation point and is exempt from pruning, so
it carries the full algebraic system: equation of state, momentum balance,
nonlinear speed and friction, node and node-source mass balances, and the
compressor-station relations. The pipe mass balance degenerates there, since
$\gamma(1-\tau^{2})\to 0$:
%
\begin{equation}
  \underbrace{\gamma\bigl(1-\tau^{2}\bigr)\,\dot\rho_{p,v}\,V_{\mathrm{cell}}}_{\;\to\;0}
  \;=\; w_v - w_{v+1} - c_v
  \qquad\Longrightarrow\qquad
  \boxed{\;w_v - w_{v+1} - c_v \;=\; 0\;}
  \label{eq:steady}
\end{equation}
%
i.e.\ the terminal point is required to be a genuine steady state. Together
with \eqref{eq:terminal} this fixes both \emph{where} the tail ends and
\emph{that} it has stopped moving.

\begin{quote}\small
\textbf{Remark (\textsc{legendre}).} Under a Legendre mesh $\tau=1$ is not a
collocation point. Only the differential state survives there, through its
continuity (extrapolation) equation
$\rho(1)=\sum_k \ell_k(1)\,\rho(\tau_k)$, and \eqref{eq:steady} cannot be
written because no flow variable exists at $\tau=1$. Pinning an
\emph{algebraic} variable such as $p_{p,v}(1)$ is then vacuous, and restoring
the algebraic equations at $\tau=1$ conflicts with the extrapolation: the
$25$ densities are extrapolated independently, one per pipe volume, whereas
the algebraic system couples them (e.g.\ three pipes share
$\texttt{node\_p}[\mathrm{sink}_1]$).
\end{quote}

\subsection{Why $\Phi_P$ is ill-posed}

$\ell_V$ and $\ell_\beta$ vanish as the tail approaches the reference, so
their integrals converge and \eqref{eq:quad} approximates a finite quantity.
The economic stage cost \eqref{eq:lP} does not: if $P\to P^{\css}\neq 0$, then
over $N$ periods
%
\begin{equation}
  \int_{T_f}^{T_f+NT} \ell_P \,\mathrm{d}t \;=\; N\,\bar\ell_P\,T
  \;\xrightarrow[N\to\infty]{}\; \infty .
\end{equation}
%
Each period contributes a fixed positive amount. The finite value returned by
\eqref{eq:quad} is therefore a truncation artefact, not an approximation: the
weights represent only
%
\begin{equation}
  \sum_{k\in\mathcal{K}} \frac{w_k}{\gamma} \;\approx\; 11.6\ \text{h}
  \;\approx\; 1.9\ \text{periods}
  \qquad (T=6\,\text{h}),
\end{equation}
%
and this effective horizon is set by $n_{\mathrm{fe}}$, $n_{\mathrm{cp}}$ and
$L$ rather than chosen.

\section{Effect of Introducing a Discount}

\subsection{Discounted tail cost}

Replace the tail integral by its discounted counterpart, with rate
$\alpha>0$:
%
\begin{equation}
  \Phi_P^{\alpha}
  \;=\;\int_{T_f}^{\infty} e^{-\alpha (t-T_f)}\,\ell_P(t)\,\mathrm{d}t .
\end{equation}
%
Using \eqref{eq:transform} and
$\operatorname{atanh}\tau=\tfrac12\ln\frac{1+\tau}{1-\tau}$,
%
\begin{equation}
  e^{-\alpha L \operatorname{atanh}\tau}
  \;=\;\left(\frac{1-\tau}{1+\tau}\right)^{p},
  \qquad p \;\equiv\; \frac{\alpha L}{2},
\end{equation}
%
so that
%
\begin{equation}
  \Phi_P^{\alpha}
  \;=\; L\int_0^1 \ell_P(\tau)\,
        \frac{(1-\tau)^{\,p-1}}{(1+\tau)^{\,p+1}}\,\mathrm{d}\tau .
  \label{eq:disc-integral}
\end{equation}

\subsection{What changes}

\paragraph{(i) The integral converges.}
The integrand in \eqref{eq:disc-integral} behaves as $(1-\tau)^{p-1}$ near
$\tau=1$, which is integrable iff
%
\begin{equation}
  p-1 > -1
  \quad\Longleftrightarrow\quad
  p>0
  \quad\Longleftrightarrow\quad
  \alpha>0 .
\end{equation}
%
The undiscounted case $p=0$ is exactly the non-integrable borderline; any
positive discount crosses it.

\paragraph{(ii) $\tau=1$ becomes an admissible quadrature node.}
The weights become
%
\begin{equation}
  w_k \;=\; \int_{\mathrm{elem}} L_k(\tau)\,
            \frac{(1-\tau)^{\,p-1}}{(1+\tau)^{\,p+1}}\,\mathrm{d}\tau ,
  \label{eq:disc-weights}
\end{equation}
%
which is finite for \emph{every} node, including $\tau=1$ (there $L_{\rm
end}$ is bounded and the singularity is integrable). The exclusion
$\tau=1\notin\mathcal{K}$ in \eqref{eq:quad} can be dropped, together with the
$L_k(1)=0$ cancellation used for the Radau branch. If $p\ge 1$ the integrand
vanishes at $\tau=1$ altogether.

\paragraph{(iii) The stage cost is unchanged.}
Equation \eqref{eq:lP} is untouched: the objective remains absolute
compressor energy, with no $\css$ reference. Only the weights
\eqref{eq:weights} are replaced by \eqref{eq:disc-weights}. This distinguishes
discounting from the rotated cost $\ell_P-\ell_P^{\css}$, which restores
convergence only by introducing a tracking reference.

\paragraph{(iv) The terminal point is no longer free of charge.}
Because $w_{\rm end}$ is finite and large for small $p$, the power at
$\tau=1$ enters the objective. In the undiscounted formulation $\tau=1$ is
constrained by \eqref{eq:steady} but never priced, so operating there costs
nothing.

\paragraph{(v) The effective horizon becomes an explicit modelling choice.}
The implicit $\approx 1.9$ periods is replaced by the discount time constant
$1/\alpha$, which carries an economic interpretation (net present value)
rather than being a by-product of $n_{\mathrm{fe}}$, $n_{\mathrm{cp}}$ and
$L$.

\subsection{What does \emph{not} change}

\begin{enumerate}
\item \textbf{The optimal policy still favours not compressing.} Minimising
      discounted energy rewards reducing $P$ exactly as before; the tail can
      still serve demand from line pack instead of from the source. That
      degeneracy lives in the constraints, not in the objective, and no
      reweighting removes it.
\item \textbf{The terminal condition \eqref{eq:terminal} and the steady-state
      requirement \eqref{eq:steady} are unaffected.}
\item \textbf{The Lyapunov terms \eqref{eq:lyap} are unaffected}, since
      $\Phi_V$ and $\Phi_\beta$ already converge without discounting.
\end{enumerate}

\subsection{Cost of the change}

A new parameter $\alpha$ enters, and the solution is no longer
discount-neutral: near-term operation is weighted more heavily than the far
tail. With $p=\alpha L/2$, the choice $p\ge 1$ (i.e.\ $\alpha\ge 2/L$)
gives the cleanest numerics but implies a short discount horizon; a longer,
more physical horizon such as $1/\alpha = 24\,$h gives $p\approx 0.21$, which
is integrable but leaves the weights strongly peaked near $\tau=1$.

\end{document}
